% Generated by tools/make_body.py; edit generator, not this file.
% UNPROOFREAD OCR TRANSCRIPTION: page-image verification remains required.
\chapter{Complex Numbers}\label{chap:1}

% Original PDF physical page 17
\sourcepage{17}
\setcounter{section}{0}\section{The Algebra Of Complex Numbers}\label{sec:1.1}
It is fundamental that real and complex numbers obey the same basic laws of arithmetic. We begin our study of complex function theory by stressing and implementing this analogy.

\setcounter{section}{1}\setcounter{subsection}{0}\subsection{Arithmetic Operations}\label{sec:1.1.1}
From elementary algebra the reader is acquainted with the imaginary unit i with the property $i^{2} = -1$. If the imaginary unit is combined with two real numbers\index{Real number} $\alpha$, $\beta$ by the processes of addition and multiplication, we obtain a complex number $\alpha + i\beta$. $\alpha$ and $\beta$ are the real and imaginary part\index{Imaginary part} of the complex number. If $\alpha = 0$, the number is said to be purely imaginary; if $\beta = 0$, it is of course real. Zero is the only number which is at once real and purely imaginary. Two complex numbers are equal if and only if they have the same real part\index{Real part} and the same imaginary part.

Addition and multiplication do not lead out from the system of complex numbers. Assuming that the ordinary rules of arithmetic apply to complex numbers we find indeed
\begin{equation*}
(\alpha + i \beta) + (\gamma + i \delta) = (\alpha + \gamma) + i (\beta + \delta)\tag{1}\label{eq:1.1}
\end{equation*}
and
\begin{equation*}
(\alpha + i \beta) (\gamma + i \delta) = (\alpha \gamma - \beta \delta) + i (\alpha \delta + \beta \gamma).\tag{2}\label{eq:1.2}
\end{equation*}
In the second identity we have made use of the relation $i^{2} = -1$. It is less obvious that division is also possible. We wish to
% Original PDF physical page 18
\sourcepage{18}
show that $(\alpha + i\beta)/(\gamma + i\delta)$ is a complex number, provided that $\gamma + i\delta \neq 0$. If the quotient is denoted by $x + iy$, we must have
\begin{equation*}
\alpha + i \beta = (\gamma + i \delta) (x + i y).
\end{equation*}
By \hyperref[eq:1.2]{(2)} this condition can be written
\begin{equation*}
\alpha + i \beta = (\gamma x - \delta y) + i (\delta x + \gamma y),
\end{equation*}
and we obtain the two equations
\begin{equation*}
\begin{array}{l} \alpha = \gamma x - \delta y \\ \beta = \delta x + \gamma y. \end{array}
\end{equation*}
This system of simultaneous linear equations has the unique solution
\begin{equation*}
x = \frac {\alpha \gamma + \beta \delta}{\gamma^ {2} + \delta^ {2}}
\end{equation*}
\begin{equation*}
y = \frac {\beta \gamma - \alpha \delta}{\gamma^ {2} + \delta^ {2}},
\end{equation*}
for we know that $\gamma^2 +\delta^2$ is not zero. We have thus the result
\begin{equation*}
\frac {\alpha + i \beta}{\gamma + i \delta} = \frac {\alpha \gamma + \beta \delta}{\gamma^ {2} + \delta^ {2}} + i \frac {\beta \gamma - \alpha \delta}{\gamma^ {2} + \delta^ {2}}.\tag{3}\label{eq:1.3}
\end{equation*}
Once the existence of the quotient has been proved, its value can be found in a simpler way. If numerator and denominator are multiplied with $\gamma - i\delta$, we find at once
\begin{equation*}
\frac {\alpha + i \beta}{\gamma + i \delta} = \frac {(\alpha + i \beta) (\gamma - i \delta)}{(\gamma + i \delta) (\gamma - i \delta)} = \frac {(\alpha \gamma + \beta \delta) + i (\beta \gamma - \alpha \delta)}{\gamma^ {2} + \delta^ {2}}.
\end{equation*}
As a special case the reciprocal of a complex number $\neq 0$ is given by
\begin{equation*}
\frac {1}{\alpha + i \beta} = \frac {\alpha - i \beta}{\alpha^ {2} + \beta^ {2}}.
\end{equation*}
We note that $i^{n}$ has only four possible values: 1, i, -1, -i. They correspond to values of n which divided by 4 leave the remainders 0, 1, 2, 3.

\subsection*{EXERCISES}
\originalnumber{1}{exercise:1.1.1.1.1}
1. Find the values of
\begin{equation*}
(1 + 2 i) ^ {3}, \quad \frac {5}{- 3 + 4 i}, \quad \left(\frac {2 + i}{3 - 2 i}\right) ^ {2}, \quad (1 + i) ^ {n} + (1 - i) ^ {n}.
\end{equation*}
% Original PDF physical page 19
\sourcepage{19}
\originalnumber{2}{exercise:1.1.1.1.2}
2. If $z = x + iy$ (x and y real), find the real and imaginary parts of
\begin{equation*}
z ^ {4}, \quad \frac {1}{z}, \quad \frac {z - 1}{z + 1}, \quad \frac {1}{z ^ {2}}.
\end{equation*}
\originalnumber{3}{exercise:1.1.1.1.3}
3. Show that
\begin{equation*}
\left(\frac {- 1 \pm i \sqrt {3}}{2}\right) ^ {3} = 1 \quad \text { and } \quad \left(\frac {\pm 1 \pm i \sqrt {3}}{2}\right) ^ {6} = 1
\end{equation*}
for all combinations of signs.

\setcounter{section}{1}\setcounter{subsection}{1}\subsection{Square Roots}\label{sec:1.1.2}
We shall now show that the square root\index{Square root} of a complex number can be found explicitly. If the given number is $\alpha + i\beta$, we are looking for a number $x + iy$ such that
\begin{equation*}
(x + i y) ^ {2} = \alpha + i \beta .
\end{equation*}
This is equivalent to the system of equations
\begin{equation*}
\begin{array}{r} x ^ {2} - y ^ {2} = \alpha \\ 2 x y = \beta . \end{array}\tag{4}\label{eq:1.4}
\end{equation*}
From these equations we obtain
\begin{equation*}
(x ^ {2} + y ^ {2}) ^ {2} = (x ^ {2} - y ^ {2}) ^ {2} + 4 x ^ {2} y ^ {2} = \alpha^ {2} + \beta^ {2}.
\end{equation*}
Hence we must have
\begin{equation*}
x ^ {2} + y ^ {2} = \sqrt {\alpha^ {2} + \beta^ {2}},
\end{equation*}
where the square root is positive or zero. Together with the first equation \hyperref[eq:1.4]{(4)} we find
\begin{equation*}
\begin{array}{l} x ^ {2} = \frac {1}{2} (\alpha + \sqrt {\alpha^ {2} + \beta^ {2}}) \\ y ^ {2} = \frac {1}{2} (- \alpha + \sqrt {\alpha^ {2} + \beta^ {2}}). \end{array}\tag{5}\label{eq:1.5}
\end{equation*}
Observe that these quantities are positive or zero regardless of the sign of $\alpha$.

The equations \hyperref[eq:1.5]{(5)} yield, in general, two opposite values for x and two for y. But these values cannot be combined arbitrarily, for the second equation \hyperref[eq:1.4]{(4)} is not a consequence of \hyperref[eq:1.5]{(5)}. We must therefore be careful to select x and y so that their product has the sign of $\beta$. This leads to the general solution
\begin{equation*}
\sqrt {\alpha + i \beta} = \pm \left(\sqrt {\frac {\alpha + \sqrt {\alpha^ {2} + \beta^ {2}}}{2}} + i \frac {\beta}{| \beta |} \sqrt {\frac {- \alpha + \sqrt {\alpha^ {2} + \beta^ {2}}}{2}}\right)\tag{6}\label{eq:1.6}
\end{equation*}
provided that $\beta \neq 0$. For $\beta = 0$ the values are $\pm \sqrt{\alpha}$ if $\alpha \geq 0$, $\pm i\sqrt{-\alpha}$
% Original PDF physical page 20
\sourcepage{20}
if\index{Associative law|(}\index{Commutative law|(} $\alpha < 0$. It is understood that all square roots of positive numbers are taken with the positive sign.

We have found that the square root of any complex number exists and has two opposite values. They coincide only if $\alpha + i\beta = 0$. They are real if $\beta = 0$, $\alpha \geq 0$ and purely imaginary if $\beta = 0$, $\alpha \leq 0$. In other words, except for zero, only positive numbers have real square roots and only negative numbers have purely imaginary square roots.

Since both square roots are in general complex, it is not possible to distinguish between the positive and negative square root of a complex number. We could of course distinguish between the upper and lower sign in \hyperref[eq:1.6]{(6)}, but this distinction is artificial and should be avoided. The correct way is to treat both square roots in a symmetric manner.

\subsection*{EXERCISES}
\originalnumber{1}{exercise:1.1.2.2.1}
1. Compute
\begin{equation*}
\sqrt {i}, \quad \sqrt {- i}, \quad \sqrt {1 + i}, \quad \sqrt {\frac {1 - i \sqrt {3}}{2}}.
\end{equation*}
\originalnumber{2}{exercise:1.1.2.2.2}
2. Find the four values of $\sqrt[4]{-1}$.

\originalnumber{3}{exercise:1.1.2.2.3}
3. Compute $\sqrt[4]{i}$ and $\sqrt[4]{-i}$.

\originalnumber{4}{exercise:1.1.2.2.4}
4. Solve the quadratic equation
\begin{equation*}
z ^ {2} + (\alpha + i \beta) z + \gamma + i \delta = 0.
\end{equation*}
\setcounter{section}{1}\setcounter{subsection}{2}\subsection{Justification}\label{sec:1.1.3}
So far our approach to complex numbers has been completely uncritical. We have not questioned the existence of a number system in which the equation $x^{2} + 1 = 0$ has a solution while all the rules of arithmetic remain in force.

We begin by recalling the characteristic properties of the real-number system which we denote by R. In the first place, R is a field\index{Field}. This means that addition and multiplication are defined, satisfying the associative, commutative, and distributive laws\index{Distributive law}. The numbers 0 and 1 are neutral elements under addition and multiplication, respectively: $\alpha + 0 = \alpha$, $\alpha \cdot 1 = \alpha$ for all $\alpha$. Moreover, the equation of subtraction $\beta + x = \alpha$ has always a solution, and the equation of division $\beta x = \alpha$ has a solution whenever $\beta \neq 0$.\hyperlink{fn:20:dagger}{\textsuperscript{\ensuremath{\dagger}}}

One shows by elementary reasoning that the neutral elements and the results of subtraction and division are unique. Also, every field is an integral domain\index{Integral domain}: $\alpha\beta = 0$ if and only if $\alpha = 0$ or $\beta = 0$.

\begingroup\renewcommand{\thefootnote}{}\footnotetext{\hypertarget{fn:20:dagger}{}\textsuperscript{\ensuremath{\dagger}} We assume that the reader has a working knowledge of elementary algebra. Although the above characterization of a field is complete, it obviously does not convey much to a student who is not already at least vaguely familiar with the concept.}\endgroup

% Original PDF physical page 21
\index{Associative law|)}\index{Commutative law|)}\sourcepage{21}
These properties are common to all fields. In addition, the field $\mathbf{R}$ has an order relation\index{Order relation} $\alpha < \beta$ (or $\beta > \alpha$). It is most easily defined in terms of the set $\mathbf{R}^{+}$ of positive real numbers: $\alpha < \beta$ if and only if $\beta - \alpha \in \mathbf{R}^{+}$. The set $\mathbf{R}^{+}$ is characterized by the following properties: (1) 0 is not a positive number; (2) if $\alpha \neq 0$ either $\alpha$ or $-\alpha$ is positive; (3) the sum and the product of two positive numbers are positive. From these conditions one derives all the usual rules for manipulation of inequalities. In particular one finds that every square $\alpha^{2}$ is either positive or zero; therefore $1 = 1^{2}$ is a positive number.

By virtue of the order relation the sums 1, 1 + 1, 1 + 1 + 1, . . . are all different. Hence R contains the natural numbers, and since it is a field it must contain the subfield formed by all rational numbers.

Finally, R satisfies the following completeness condition: every increasing and bounded sequence of real numbers has a limit. Let $\alpha_{1} < \alpha_{2} < \alpha_{3} < \cdots < \alpha_{n} < \cdots$, and assume the existence of a real number B such that $\alpha_{n} < B$ for all n. Then the completeness condition requires the existence of a number $A = \lim_{n \to \infty} \alpha_{n}$ with the following property: given any $\varepsilon > 0$ there exists a natural number $n_{0}$ such that $A - \varepsilon < \alpha_{n} < A$ for all $n > n_{0}$.

Our discussion of the real-number system is incomplete inasmuch as we have not proved the existence and uniqueness (up to isomorphisms) of a system R with the postulated properties.\hyperlink{fn:21:dagger}{\textsuperscript{\ensuremath{\dagger}}} The student who is not thoroughly familiar with one of the constructive processes by which real numbers can be introduced should not fail to fill this gap by consulting any textbook in which a full axiomatic treatment of real numbers is given.

The equation $x^{2} + 1 = 0$ has no solution in R, for $\alpha^{2} + 1$ is always positive. Suppose now that a field F can be found which contains R as a subfield, and in which the equation $x^{2} + 1 = 0$ can be solved. Denote a solution by i. Then $x^{2} + 1 = (x + i)(x - i)$, and the equation $x^{2} + 1 = 0$ has exactly two roots in F, i and -i. Let C be the subset of F consisting of all elements which can be expressed in the form $\alpha + i\beta$ with real $\alpha$ and $\beta$. This representation is unique, for $\alpha + i\beta = \alpha' + i\beta'$ implies $\alpha - \alpha' = -i(\beta - \beta')$; hence $(\alpha - \alpha')^{2} = -(\beta - \beta')^{2}$, and this is possible only if $\alpha = \alpha'$, $\beta = \beta'$.

The subset C is a subfield of F. In fact, except for trivial verifications which the reader is asked to carry out, this is exactly what was shown in Sec.~\hyperref[sec:1.1.1]{1.1}. What is more, the structure of C is independent of F. For if $F'$ is another field containing R and a root $i'$ of the equation $x^{2} + 1 = 0$,
\begingroup\renewcommand{\thefootnote}{}\footnotetext{\hypertarget{fn:21:dagger}{}\textsuperscript{\ensuremath{\dagger}} An isomorphism\index{Isomorphism} between two fields is a one-to-one correspondence which preserves sums and products. The word is used quite generally to indicate a correspondence which is one to one and preserves all relations that are considered important in a given connection.}\endgroup
% Original PDF physical page 22
\sourcepage{22}
the\index{Conjugate number|(} corresponding subset $C'$ is formed by all elements $\alpha + i'\beta$. There is a one-to-one correspondence between C and $C'$ which associates $\alpha + i\beta$ and $\alpha + i'\beta$, and this correspondence is evidently a field isomorphism. It is thus demonstrated that C and $C'$ are isomorphic.

We now define the field of complex numbers to be the subfield C of an arbitrarily given F. We have just seen that the choice of F makes no difference, but we have not yet shown that there exists a field F with the required properties. In order to give our definition a meaning it remains to exhibit a field F which contains R (or a subfield isomorphic with R) and in which the equation $x^{2} + 1 = 0$ has a root.

There are many ways in which such a field can be constructed. The simplest and most direct method is the following: Consider all expressions of the form $\alpha + i\beta$ where $\alpha, \beta$ are real numbers while the signs + and i are pure symbols (+ does not indicate addition, and i is not an element of a field). These expressions are elements of a field F in which addition and multiplication are defined by \hyperref[eq:1.1]{(1)} and \hyperref[eq:1.2]{(2)} (observe the two different meanings of the sign +). The elements of the particular form $\alpha + i0$ are seen to constitute a subfield isomorphic to R, and the element $0 + i1$ satisfies the equation $x^{2} + 1 = 0$; we obtain in fact $(0 + i1)^{2} = -(1 + i0)$. The field F has thus the required properties; moreover, it is identical with the corresponding subfield C, for we can write
\begin{equation*}
\alpha + i \beta = (\alpha + i 0) + \beta (0 + i 1).
\end{equation*}
The existence of the complex-number field is now proved, and we can go back to the simpler notation $\alpha + i\beta$ where the + indicates addition in C and i is a root of the equation $x^{2} + 1 = 0$.

EXERCISES (For students with a background in algebra)

1. Show that the system of all matrices of the special form
\begin{equation*}
\left( \begin{array}{c c} \alpha & \beta \\ - \beta & \alpha \end{array} \right),
\end{equation*}
combined by matrix addition and matrix multiplication, is isomorphic to the field of complex numbers.

2. Show that the complex-number system can be thought of as the field of all polynomials with real coefficients modulo the irreducible polynomial $x^{2} + 1$.

\setcounter{section}{1}\setcounter{subsection}{3}\subsection{Conjugation, Absolute Value}\index{Absolute value|(}\index{Absolute value}\label{sec:1.1.4}
A complex number can be denoted either by a single letter $a$, representing an element of the field $\mathbf{C}$, or in the form $\alpha + i\beta$ with real $\alpha$ and $\beta$. Other standard notations are $z = x + iy$, $\zeta = \xi + i\eta$, $w = u + iv$, and when used in this connection it
% Original PDF physical page 23
\sourcepage{23}
is tacitly understood that $x, y, \xi, \eta, u, v$ are real numbers. The real and imaginary part of a complex number $a$ will also be denoted by $\operatorname{Re} a$, $\operatorname{Im} a$.

In deriving the rules for complex addition and multiplication we used only the fact that $i^{2} = -1$. Since -i has the same property, all rules must remain valid if i is everywhere replaced by -i. Direct verification shows that this is indeed so. The transformation which replaces $\alpha + i\beta$ by $\alpha - i\beta$ is called complex conjugation, and $\alpha - i\beta$ is the conjugate of $\alpha + i\beta$. The conjugate of a is denoted by $\bar{a}$. A number is real if and only if it is equal to its conjugate. The conjugation is an involutory transformation\index{Involutory transformation}: this means that $\overline{\bar{a}} = a$.

The formulas
\begin{equation*}
\operatorname{Re} a = \frac {a + \bar {a}}{2}, \quad \operatorname{Im} a = \frac {a - \bar {a}}{2 i}
\end{equation*}
express the real and imaginary part in terms of the complex number and its conjugate. By systematic use of the notations a and $\bar{a}$ it is hence possible to dispense with the use of separate letters for the real and imaginary part. It is more convenient, though, to make free use of both notations.

The fundamental property of conjugation is the one already referred to, namely, that
\begin{equation*}
\begin{array}{r} \overline {{a + b}} = \bar {a} + \bar {b} \\ \overline {{a b}} = \bar {a} \cdot \bar {b}. \end{array}
\end{equation*}
The corresponding property for quotients is a consequence: if $ax = b$, then $\bar{a}\bar{x} = \bar{b}$, and hence $(\overline{b / a}) = \bar{b} / \bar{a}$. More generally, let $R(a,b,c,\ldots)$ stand for any rational operation applied to the complex numbers $a, b, c, \ldots$. Then
\begin{equation*}
\overline {{R (a , b , c , \dots)}} = R (\bar {a}, \bar {b}, \bar {c}, \dots).
\end{equation*}
As an application, consider the equation
\begin{equation*}
c _ {0} z ^ {n} + c _ {1} z ^ {n - 1} + \dots + c _ {n - 1} z + c _ {n} = 0.
\end{equation*}
If $\zeta$ is a root of this equation, then $\bar{\zeta}$ is a root of the equation
\begin{equation*}
\bar {c} _ {0} z ^ {n} + \bar {c} _ {1} z ^ {n - 1} + \cdot \cdot \cdot + \bar {c} _ {n - 1} z + \bar {c} _ {n} = 0.
\end{equation*}
In particular, if the coefficients are real, $\zeta$ and $\bar{\zeta}$ are roots of the same equation, and we have the familiar theorem that the nonreal roots of an equation with real coefficients occur in pairs of conjugate roots.

The product $a\bar{a} = \alpha^{2} + \beta^{2}$ is always positive or zero. Its nonnegative square root is called the modulus\index{Modulus} or absolute value\index{Absolute value} of the complex number a; it is denoted by $|a|$. The terminology and notation are justified by
% Original PDF physical page 24
\sourcepage{24}
the fact that the modulus of a real number coincides with its numerical value taken with the positive sign.

We repeat the definition
\begin{equation*}
a \bar {a} = | a | ^ {2},
\end{equation*}
where $|a| \geq 0$, and observe that $|\bar{a}| = |a|$. For the absolute value\index{Absolute value} of a product we obtain
\begin{equation*}
\left| a b \right| ^ {2} = a b \cdot \overline {{{a b}}} = a b \bar {a} \bar {b} = a \bar {a} b \bar {b} = | a | ^ {2} | b | ^ {2},
\end{equation*}
and hence
\begin{equation*}
\vert a b \vert = \vert a \vert \cdot \vert b \vert
\end{equation*}
since both are $\geq 0$. In words:

The absolute value of a product is equal to the product of the absolute values of the factors.

It is clear that this property extends to arbitrary finite products:
\begin{equation*}
\left| a _ {1} a _ {2} \cdot \cdot \cdot a _ {n} \right| = \left| a _ {1} \right| \cdot \left| a _ {2} \right| \cdot \cdot \cdot \left| a _ {n} \right|.
\end{equation*}
The quotient $a / b$, $b \neq 0$, satisfies $b(a / b) = a$, and hence we have also $|b| \cdot |a / b| = |a|$, or
\begin{equation*}
\left| \frac {a}{b} \right| = \frac {| a |}{| b |}.
\end{equation*}
The formula for the absolute value of a sum is not as simple. We find
\begin{equation*}
\vert a + b \vert^ {2} = (a + b) (\bar {a} + \bar {b}) = a \bar {a} + (a \bar {b} + b \bar {a}) + b \bar {b}
\end{equation*}
or
\begin{equation*}
\vert a + b \vert^ {2} = \vert a \vert^ {2} + \vert b \vert^ {2} + 2 \operatorname{Re} a \bar {b}.\tag{7}\label{eq:1.7}
\end{equation*}
The corresponding formula for the difference is
\begin{equation*}
\vert a - b \vert^ {2} = \vert a \vert^ {2} + \vert b \vert^ {2} - 2 \operatorname{Re} a \bar {b},\tag{7'}\label{eq:1.7'}
\end{equation*}
and by addition we obtain the identity
\begin{equation*}
\vert a + b \vert^ {2} + \vert a - b \vert^ {2} = 2 (\vert a \vert^ {2} + \vert b \vert^ {2}).\tag{8}\label{eq:1.8}
\end{equation*}
\subsection*{EXERCISES}
\originalnumber{1}{exercise:1.1.4.3.1}
1. Verify by calculation that the values of
\begin{equation*}
\frac {z}{z ^ {2} + 1}
\end{equation*}
for $z = x + iy$ and $z = x - iy$ are conjugate.

\originalnumber{2}{exercise:1.1.4.3.2}
2. Find the absolute values of
\begin{equation*}
- 2 i (3 + i) (2 + 4 i) (1 + i) \quad \text { and } \quad \frac {(3 + 4 i) (- 1 + 2 i)}{(- 1 - i) (3 - i)}.
\end{equation*}
% Original PDF physical page 25
\index{Absolute value|)}\index{Conjugate number|)}\sourcepage{25}
\originalnumber{3}{exercise:1.1.4.3.3}
3. Prove that
\begin{equation*}
\left| \frac {a - b}{1 - \bar {a} b} \right| = 1
\end{equation*}
if either $|a| = 1$ or $|b| = 1$. What exception must be made if $|a| = |b| = 1$?

\originalnumber{4}{exercise:1.1.4.3.4}
4. Find the conditions under which the equation $az + b\bar{z} + c = 0$ in one complex unknown has exactly one solution, and compute that solution.

\originalnumber{5}{exercise:1.1.4.3.5}
5. Prove Lagrange's identity\index{Identity, Lagrange's}\index{Lagrange's identity} in the complex form
\begin{equation*}
\left| \sum_ {i = 1} ^ {n} a _ {i} b _ {i} \right| ^ {2} = \sum_ {i = 1} ^ {n} | a _ {i} | ^ {2} \sum_ {i = 1} ^ {n} | b _ {i} | ^ {2} - \sum_ {1 \leq i <   j \leq n} | a _ {i} \bar {b} _ {j} - a _ {j} \bar {b} _ {i} | ^ {2}.
\end{equation*}
\setcounter{section}{1}\setcounter{subsection}{4}\subsection{Inequalities}\label{sec:1.1.5}
We shall now prove some important inequalities which will be of constant use. It is perhaps well to point out that there is no order relation in the complex-number system, and hence all inequalities must be between real numbers.

From the definition of the absolute value we deduce the inequalities
\begin{equation*}
\begin{array}{l} - | a | \leq \text {   Re   } a \leq | a | \\ - | a | \leq \text {   Im   } a \leq | a |. \end{array}\tag{9}\label{eq:1.9}
\end{equation*}
The equality $\operatorname{Re} a = |a|$ holds if and only if $a$ is real and $\geq 0$.

If \hyperref[eq:1.9]{(9)} is applied to \hyperref[eq:1.7]{(7)}, we obtain
\begin{equation*}
\vert a + b \vert^ {2} \leq (\vert a \vert + \vert b \vert) ^ {2}
\end{equation*}
and hence
\begin{equation*}
| a + b | \leq | a | + | b |.\tag{10}\label{eq:1.10}
\end{equation*}
This is called the triangle inequality\index{Triangle inequality} for reasons which will emerge later. By induction it can be extended to arbitrary sums:
\begin{equation*}
\left| a _ {1} + a _ {2} + \dots + a _ {n} \right| \leq \left| a _ {1} \right| + \left| a _ {2} \right| + \dots + \left| a _ {n} \right|.\tag{11}\label{eq:1.11}
\end{equation*}
The absolute value of a sum is at most equal to the sum of the absolute values of the terms.

The reader is well aware of the importance of the estimate \hyperref[eq:1.11]{(11)} in the real case, and we shall find it no less important in the theory of complex numbers.

Let us determine all cases of equality in \hyperref[eq:1.11]{(11)}. In \hyperref[eq:1.10]{(10)} the equality holds if and only if $a\bar{b} \geq 0$ (it is convenient to let $c > 0$ indicate that $c$ is real and positive). If $b \neq 0$ this condition can be written in the form $|b|^2(a/b) \geq 0$, and it is hence equivalent to $a/b \geq 0$. In the general
% Original PDF physical page 26
\sourcepage{26}
case we proceed as follows: Suppose that equality holds in \hyperref[eq:1.11]{(11)}; then
\begin{equation*}
\begin{array}{c} | a _ {1} | + | a _ {2} | + \dots + | a _ {n} | = | (a _ {1} + a _ {2}) + a _ {3} + \dots + a _ {n} | \\ \leq | a _ {1} + a _ {2} | + | a _ {3} | + \dots + | a _ {n} | \leq | a _ {1} | + | a _ {2} | + \dots + | a _ {n} |. \end{array}
\end{equation*}
Hence $|a_1 + a_2| = |a_1| + |a_2|$, and if $a_2 \neq 0$ we conclude that $a_1 / a_2 \geq 0$. But the numbering of the terms is arbitrary; thus the ratio of any two nonzero terms must be positive. Suppose conversely that this condition is fulfilled. Assuming that $a_1 \neq 0$ we obtain
\begin{equation*}
\begin{aligned}
|a_1+a_2+\dots+a_n|&=|a_1|\left|1+\frac{a_2}{a_1}+\dots+\frac{a_n}{a_1}\right|\\
&=|a_1|\left(1+\frac{a_2}{a_1}+\dots+\frac{a_n}{a_1}\right)\\
&=|a_1|\left(1+\frac{|a_2|}{|a_1|}+\dots+\frac{|a_n|}{|a_1|}\right)\\
&=|a_1|+|a_2|+\dots+|a_n|.
\end{aligned}
\end{equation*}
To sum up: the sign of equality holds in \hyperref[eq:1.11]{(11)} if and only if the ratio of any two nonzero terms is positive.

By \hyperref[eq:1.10]{(10)} we have also
\begin{equation*}
| a | = | (a - b) + b | \leq | a - b | + | b |
\end{equation*}
or
\begin{equation*}
| a | - | b | \leq | a - b |.
\end{equation*}
For the same reason $|b| - |a| \leq |a - b|$, and these inequalities can be combined to
\begin{equation*}
| a - b | \geq | | a | - | b | |.\tag{12}\label{eq:1.12}
\end{equation*}
Of course the same estimate can be applied to $|a + b|$.

A special case of \hyperref[eq:1.10]{(10)} is the inequality
\begin{equation*}
| \alpha + i \beta | \leq | \alpha | + | \beta |\tag{13}\label{eq:1.13}
\end{equation*}
which expresses that the absolute value of a complex number is at most equal to the sum of the absolute values of the real and imaginary part.

Many other inequalities whose proof is less immediate are also of frequent use. Foremost is Cauchy's inequality\index{Cauchy's inequality} which states that
\begin{equation*}
\left| a _ {1} b _ {1} + \dots + a _ {n} b _ {n} \right| ^ {2} \leq \left(\left| a _ {1} \right| ^ {2} + \dots + \left| a _ {n} \right| ^ {2}\right) \left(\left| b _ {1} \right| ^ {2} + \dots + \left| b _ {n} \right| ^ {2}\right)
\end{equation*}
or, in shorter notation,
\begin{equation*}
\left| \sum_ {i = 1} ^ {n} a _ {i} b _ {i} \right| ^ {2} \leq \sum_ {i = 1} ^ {n} | a _ {i} | ^ {2} \sum_ {i = 1} ^ {n} | b _ {i} | ^ {2}. \text{\hyperlink{fn:26:dagger}{\textsuperscript{\ensuremath{\dagger}}}}\tag{14}\label{eq:1.14}
\end{equation*}
\begingroup\renewcommand{\thefootnote}{}\footnotetext{\hypertarget{fn:26:dagger}{}\textsuperscript{\ensuremath{\dagger}} $i$ is a convenient summation index and, used as a subscript, cannot be confused with the imaginary unit. It seems pointless to ban its use.}\endgroup

% Original PDF physical page 27
\sourcepage{27}
To prove it, let $\lambda$ denote an arbitrary complex number. We obtain by \hyperref[eq:1.7]{(7)}
\begin{equation*}
\sum_ {i = 1} ^ {n} | a _ {i} - \lambda \bar {b} _ {i} | ^ {2} = \sum_ {i = 1} ^ {n} | a _ {i} | ^ {2} + | \lambda | ^ {2} \sum_ {i = 1} ^ {n} | b _ {i} | ^ {2} - 2 \operatorname{Re} \bar {\lambda} \sum_ {i = 1} ^ {n} a _ {i} b _ {i}.\tag{15}\label{eq:1.15}
\end{equation*}
This expression is $\geq 0$ for all $\lambda$. We can choose
\begin{equation*}
\lambda = \frac {\sum_ {1} ^ {n} a _ {i} b _ {i}}{\sum_ {1} ^ {n} | b _ {i} | ^ {2}},
\end{equation*}
for if the denominator should vanish there is nothing to prove. This choice is not arbitrary, but it is dictated by the desire to make the expression \hyperref[eq:1.15]{(15)} as small as possible. Substituting in \hyperref[eq:1.15]{(15)} we find, after simplifications,
\begin{equation*}
\sum_ {1} ^ {n} \left| a _ {i} \right| ^ {2} - \frac {\left| \sum_ {1} ^ {n} a _ {i} b _ {i} \right| ^ {2}}{\sum_ {1} ^ {n} \left| b _ {i} \right| ^ {2}} \geq 0
\end{equation*}
which is equivalent to \hyperref[eq:1.14]{(14)}.

From \hyperref[eq:1.15]{(15)} we conclude further that the sign of equality holds in \hyperref[eq:1.14]{(14)} if and only if the $a_{i}$ are proportional to the $\bar{b}_{i}$.

Cauchy's inequality can also be proved by means of Lagrange's identity (Sec.~\hyperref[sec:1.1.4]{1.4}, Ex. \hyperref[exercise:1.1.4.3.4]{4}).

\subsection*{EXERCISES}
\originalnumber{1}{exercise:1.1.5.4.1}
1. Prove that
\begin{equation*}
\left| \frac {a - b}{1 - \bar {a} b} \right| <   1
\end{equation*}
if $|a| < 1$ and $|b| < 1$.

\originalnumber{2}{exercise:1.1.5.4.2}
2. Prove Cauchy's inequality by induction.

\originalnumber{3}{exercise:1.1.5.4.3}
3. If $|a_{i}| < 1, \lambda_{i} \geq 0$ for $i = 1, \ldots, n$ and $\lambda_{1} + \lambda_{2} + \cdots + \lambda_{n} = 1$, show that
\begin{equation*}
\left| \lambda_ {1} a _ {1} + \lambda_ {2} a _ {2} + \dots + \lambda_ {n} a _ {n} \right| <   1.
\end{equation*}
\originalnumber{4}{exercise:1.1.5.4.4}
4. Show that there are complex numbers z satisfying
\begin{equation*}
\vert z - a \vert + \vert z + a \vert = 2 \vert c \vert
\end{equation*}
% Original PDF physical page 28
\sourcepage{28}
if and only if $|a| \leq |c|$. If this condition is fulfilled, what are the smallest and largest values of $|z|$?

\setcounter{section}{1}\section{The Geometric Representation Of Complex Numbers}\label{sec:1.2}
With respect to a given rectangular coordinate system in a plane, the complex number $a = \alpha + i\beta$ can be represented by the point\index{Point} with coordinates $(\alpha, \beta)$. This representation is constantly used, and we shall often speak of the point a as a synonym of the number a. The first coordinate axis (x-axis) takes the name of real axis\index{Axis!real}, and the second coordinate axis (y-axis) is called the imaginary axis\index{Axis!imaginary}\index{Imaginary axis}. The plane itself is referred to as the complex plane\index{Plane!complex}.

The geometric representation derives its usefulness from the vivid mental pictures associated with a geometric language. We take the point of view, however, that all conclusions in analysis should be derived from the properties of real numbers, and not from the axioms of geometry. For this reason we shall use geometry only for descriptive purposes, and not for valid proof, unless the language is so thinly veiled that the analytic interpretation is self-evident. This attitude relieves us from the exigencies of rigor in connection with geometric considerations.

\setcounter{section}{2}\setcounter{subsection}{0}\subsection{Geometric Addition and Multiplication}\label{sec:1.2.1}
The addition of complex numbers can be visualized as vector\index{Vector} addition\index{Vector addition}. To this end we let a complex number be represented not only by a point, but also by a vector pointing from the origin to the point. The number, the point, and the vector will all be denoted by the same letter a. As usual we identify all vectors which can be obtained from each other by parallel displacements.

Place a second vector b so that its initial point coincides with the end point of a. Then $a + b$ is represented by the vector from the initial point of a to the end point of b. To construct the difference b - a we draw both vectors a and b from the same initial point; then b - a points from the end point of a to the end point of b. Observe that $a + b$ and a - b are the diagonals in a parallelogram with the sides a and b (Fig.~\hyperref[fig:1-1]{1-1}).

An additional advantage of the vector representation is that the length of the vector a is equal to $|a|$. Hence the distance between the points a and b is $|a - b|$. With this interpretation the triangle inequality $|a + b| \leq |a| + |b|$ and the identity $|a + b|^{2} + |a - b|^{2} = 2(|a|^{2} + |b|^{2})$ become familiar geometric theorems.

The point a and its conjugate $\bar{a}$ lie symmetrically with respect to the real axis. The symmetric point of a with respect to the imaginary axis is


% Original PDF physical page 29
\sourcepage{29}
\bookfigure{29}{1}{FIG. 1-1. Vector addition.}
$-\bar{a}$. The four points $a, -\bar{a}, -a, \bar{a}$ are the vertices of a rectangle which is symmetric with respect to both axes.

In order to derive a geometric interpretation of the product of two complex numbers we introduce polar coordinates. If the polar coordinates of the point $(\alpha,\beta)$ are $(r,\varphi)$, we know that
\begin{equation*}
\begin{array}{l} \alpha = r \cos \varphi \\ \beta = r \sin \varphi . \end{array}
\end{equation*}
Hence we can write $a = \alpha + i\beta = r(\cos \varphi + i \sin \varphi)$. In this trigonometric form of a complex number r is always $\geq 0$ and equal to the modulus $|a|$. The polar angle $\varphi$ is called the argument\index{Argument} or amplitude\index{Amplitude} of the complex number, and we denote it by arg a.

Consider two complex numbers $a_{1}=r_{1}(\cos\varphi_{1}+i\sin\varphi_{1})$ and $a_{2}=r_{2}(\cos\varphi_{2}+i\sin\varphi_{2})$. Their product can be written in the from $a_{1}a_{2}=r_{1}r_{2}[(\cos\varphi_{1}\cos\varphi_{2}-\sin\varphi_{1}\sin\varphi_{2})+i(\sin\varphi_{1}\cos\varphi_{2}+\cos\varphi_{1}\sin\varphi_{2})]$. By means of the addition theorems of the cosine and the sine this expression can be simplified to
\begin{equation*}
a _ {1} a _ {2} = r _ {1} r _ {2} \left[ \cos \left(\varphi_ {1} + \varphi_ {2}\right) + i \sin \left(\varphi_ {1} + \varphi_ {2}\right) \right].\tag{16}\label{eq:1.16}
\end{equation*}
We recognize that the product has the modulus $r_{1}r_{2}$ and the argument $\varphi_{1} + \varphi_{2}$. The latter result is new, and we express it through the equation
\begin{equation*}
\arg \left(a _ {1} a _ {2}\right) = \arg a _ {1} + \arg a _ {2}.\tag{17}\label{eq:1.17}
\end{equation*}
It is clear that this formula can be extended to arbitrary products, and we can therefore state:

The argument of a product is equal to the sum of the arguments of the factors.

This is fundamental. The rule that we have just formulated gives a deep and unexpected justification of the geometric representation of complex numbers. We must be fully aware, however, that the manner in which we have arrived at the formula \hyperref[eq:1.17]{(17)} violates our principles. In the
% Original PDF physical page 30
\sourcepage{30}
first place the equation \hyperref[eq:1.17]{(17)} is between angles rather than between numbers, and secondly its proof rested on the use of trigonometry. Thus it remains to define the argument in analytic terms and to prove \hyperref[eq:1.17]{(17)} by purely analytic means. For the moment we postpone this proof and shall be content to discuss the consequences of \hyperref[eq:1.17]{(17)} from a less critical standpoint.

We remark first that the argument of 0 is not defined, and hence \hyperref[eq:1.17]{(17)} has a meaning only if $a_1$ and $a_2$ are $\neq 0$. Secondly, the polar angle\index{Angle} is determined only up to multiples of $360^\circ$. For this reason, if we want to interpret \hyperref[eq:1.17]{(17)} numerically, we must agree that multiples of $360^\circ$ shall not count.

By means of \hyperref[eq:1.17]{(17)} a simple geometric construction of the product $a_{1}a_{2}$ can be obtained. It follows indeed that the triangle with the vertices 0, 1, $a_{1}$ is similar to the triangle whose vertices are 0, $a_{2}$, $a_{1}a_{2}$. The points 0, 1, $a_{1}$, and $a_{2}$ being given, this similarity determines the point $a_{1}a_{2}$ (Fig.~\hyperref[fig:1-2]{1-2}). In the case of division \hyperref[eq:1.17]{(17)} is replaced by
\begin{equation*}
\arg \frac {a _ {2}}{a _ {1}} = \arg a _ {2} - \arg a _ {1}.\tag{18}\label{eq:1.18}
\end{equation*}
The geometric construction is the same, except that the similar triangles are now 0, 1, $a_1$ and 0, $a_2 / a_1$, $a_2$.

Remark: A perfectly acceptable way to define angles and arguments would be to apply the familiar methods of calculus which permit us to express the length of a circular arc as a definite integral. This leads to a correct definition of the trigonometric functions, and to a computational proof of the addition theorems.

The reason we do not follow this path is that complex analysis, as
\bookfigure{30}{1}{FIG. 1-2. Vector multiplication.}
% Original PDF physical page 31
\sourcepage{31}
opposed to real analysis, offers a much more direct approach. The clue lies in a direct connection between the exponential function and the trigonometric functions, to be derived in Chap.~\hyperref[chap:2]{2}, Sec. 5. Until we reach this point the reader is asked to subdue his quest for complete rigor.

\subsection*{EXERCISES}
\originalnumber{1}{exercise:1.2.1.5.1}
1. Find the symmetric points of a with respect to the lines which bisect the angles between the coordinate axes.

\originalnumber{2}{exercise:1.2.1.5.2}
2. Prove that the points $a_{1}, a_{2}, a_{3}$ are vertices of an equilateral triangle if and only if $a_{1}^{2} + a_{2}^{2} + a_{3}^{2} = a_{1}a_{2} + a_{2}a_{3} + a_{3}a_{1}$.

\originalnumber{3}{exercise:1.2.1.5.3}
3. Suppose that $a$ and $b$ are two vertices of a square. Find the two other vertices in all possible cases.

\originalnumber{4}{exercise:1.2.1.5.4}
4. Find the center and the radius of the circle which circumscribes the triangle with vertices $a_{1}, a_{2}, a_{3}$. Express the result in symmetric form.

\setcounter{section}{2}\setcounter{subsection}{1}\subsection{The Binomial Equation}\index{Binomial equation|(}\index{Binomial equation}\label{sec:1.2.2}
From the preceding results we derive that the powers of $a = r(\cos \varphi + i\sin \varphi)$ are given by
\begin{equation*}
a ^ {n} = r ^ {n} (\cos n \varphi + i \sin n \varphi).\tag{19}\label{eq:1.19}
\end{equation*}
This formula is trivially valid for $n = 0$, and since
\begin{equation*}
a ^ {- 1} = r ^ {- 1} (\cos \varphi - i \sin \varphi) = r ^ {- 1} [ \cos (- \varphi) + i \sin (- \varphi) ]
\end{equation*}
it holds also when $n$ is a negative integer.

For $r = 1$ we obtain de Moivre's formula\index{de Moivre's formula}
\begin{equation*}
(\cos \varphi + i \sin \varphi) ^ {n} = \cos n \varphi + i \sin n \varphi\tag{20}\label{eq:1.20}
\end{equation*}
which provides an extremely simple way to express $\cos n\varphi$ and $\sin n\varphi$ in terms of $\cos \varphi$ and $\sin \varphi$.

To find the $n$th root of a complex number $a$ we have to solve the equation
\begin{equation*}
z ^ {n} = a.\tag{21}\label{eq:1.21}
\end{equation*}
Supposing that $a \neq 0$ we write $a = r(\cos \varphi + i \sin \varphi)$ and
\begin{equation*}
z = \rho (\cos \theta + i \sin \theta).
\end{equation*}
Then \hyperref[eq:1.21]{(21)} takes the form
\begin{equation*}
\rho^ {n} (\cos n \theta + i \sin n \theta) = r (\cos \varphi + i \sin \varphi).\tag{22}\label{eq:1.22}
\end{equation*}
% Original PDF physical page 32
\sourcepage{32}
This equation is certainly fulfilled if $\rho^{n}=r$ and $n\theta=\varphi$. Hence we obtain the root
\begin{equation*}
z = \sqrt [ n ]{r} \left(\cos \frac {\varphi}{n} + i \sin \frac {\varphi}{n}\right),
\end{equation*}
where $\sqrt[n]{r}$ denotes the positive nth root of the positive number r.

But this is not the only solution. In fact, \hyperref[eq:1.22]{(22)} is also fulfilled if $n\theta$ differs from $\varphi$ by a multiple of the full angle. If angles are expressed in radians the full angle is $2\pi$, and we find that \hyperref[eq:1.22]{(22)} is satisfied if and only if
\begin{equation*}
\theta = \frac {\varphi}{n} + k \cdot \frac {2 \pi}{n},
\end{equation*}
where k is any integer. However, only the values $k = 0, 1, \ldots, n - 1$ give different values of z. Hence the complete solution of the equation \hyperref[eq:1.21]{(21)} is given by
\begin{equation*}
z = \sqrt [ n ]{r} \left[ \cos \left(\frac {\varphi}{n} + k \frac {2 \pi}{n}\right) + i \sin \left(\frac {\varphi}{n} + k \frac {2 \pi}{n}\right) \right], \quad k = 0, 1, \dots , n - 1.
\end{equation*}
There are n nth roots of any complex number $\neq 0$. They have the same modulus, and their arguments are equally spaced.

Geometrically, the $n$th roots are the vertices of a regular polygon with $n$ sides.

The case $a = 1$ is particularly important. The roots of the equation $z^n = 1$ are called $n$th roots of unity, and if we set
\begin{equation*}
\omega = \cos {\frac {2 \pi}{n}} + i \sin {\frac {2 \pi}{n}}\tag{23}\label{eq:1.23}
\end{equation*}
all the roots can be expressed by 1, $\omega$, $\omega^2$, $\ldots$, $\omega^{n-1}$. It is also quite evident that if $\sqrt[n]{a}$ denotes any $n$th root of $a$, then all the $n$th roots can be expressed in the form $\omega^k \cdot \sqrt[n]{a}$, $k = 0, 1, \ldots, n - 1$.

\subsection*{EXERCISES}
\originalnumber{1}{exercise:1.2.2.6.1}
1. Express $\cos 3\varphi$, $\cos 4\varphi$, and $\sin 5\varphi$ in terms of $\cos \varphi$ and $\sin \varphi$.

\originalnumber{2}{exercise:1.2.2.6.2}
2. Simplify $1 + \cos \varphi + \cos 2\varphi + \cdots + \cos n\varphi$ and $\sin \varphi + \sin 2\varphi + \cdots + \sin n\varphi$.

\originalnumber{3}{exercise:1.2.2.6.3}
3. Express the fifth and tenth roots of unity in algebraic form.

\originalnumber{4}{exercise:1.2.2.6.4}
4. If $\omega$ is given by \hyperref[eq:1.23]{(23)}, prove that
\begin{equation*}
1 + \omega^ {h} + \omega^ {2 h} + \dots + \omega^ {(n - 1) h} = 0
\end{equation*}
for any integer h which is not a multiple of n.


% Original PDF physical page 33
\index{Binomial equation|)}\sourcepage{33}
\originalnumber{5}{exercise:1.2.2.6.5}
5. What is the value of
\begin{equation*}
1 - \omega^ {h} + \omega^ {2 h} - \dots + (- 1) ^ {n - 1} \omega^ {(n - 1) h}
\end{equation*}
\setcounter{section}{2}\setcounter{subsection}{2}\subsection{Analytic Geometry}\label{sec:1.2.3}
In classical analytic geometry\index{Analytic geometry} the equation of a locus is expressed as a relation between x and y. It can just as well be expressed in terms of z and $\bar{z}$, sometimes to distinct advantage. The thing to remember is that a complex equation is ordinarily equivalent to two real equations; in order to obtain a genuine locus these equations should be essentially the same.

For instance, the equation of a circle is $|z - a| = r$. In algebraic form it can be rewritten as $(z - a)(\bar{z} - \bar{a}) = r^{2}$. The fact that this equation is invariant under complex conjugation is an indication that it represents a single real equation.

A straight line\index{Straight line} in the complex plane can be given by a parametric equation $z = a + bt$, where a and b are complex numbers and $b \neq 0$; the parameter t runs through all real values. Two equations $z = a + bt$ and $z = a' + b't$ represent the same line if and only if $a' - a$ and $b'$ are real multiples of b. The lines are parallel whenever $b'$ is a real multiple of b, and they are equally directed if $b'$ is a positive multiple of b. The direction of a directed line can be identified with arg b. The angle between $z = a + bt$ and $z = a' + b't$ is arg $b'/b$; observe that it depends on the order in which the lines are named. The lines are orthogonal to each other if $b'/b$ is purely imaginary.

Problems of finding intersections between lines and circles, parallel or orthogonal lines, tangents, and the like usually become exceedingly simple when expressed in complex form.

An inequality $|z - a| < r$ describes the inside of a circle. Similarly, a directed line $z = a + bt$ determines a right half plane consisting of all points z with $\operatorname{Im}(z - a)/b < 0$ and a left half plane with $\operatorname{Im}(z - a)/b > 0$. An easy argument shows that this distinction is independent of the parametric representation.

\subsection*{EXERCISES}
\originalnumber{1}{exercise:1.2.3.7.1}
1. When does $az + b\bar{z} + c = 0$ represent a line?

\originalnumber{2}{exercise:1.2.3.7.2}
2. Write the equation of an ellipse, hyperbola, parabola in complex form.

\originalnumber{3}{exercise:1.2.3.7.3}
3. Prove that the diagonals of a parallelogram bisect each other and that the diagonals of a rhombus are orthogonal.

\originalnumber{4}{exercise:1.2.3.7.4}
4. Prove analytically that the midpoints of parallel chords to a circle lie on a diameter perpendicular to the chords.

\originalnumber{5}{exercise:1.2.3.7.5}
5. Show that all circles that pass through a and $1/\bar{a}$ intersect the circle $|z| = 1$ at right angles.


% Original PDF physical page 34
\sourcepage{34}
\setcounter{section}{2}\setcounter{subsection}{3}\subsection{The Spherical Representation}\label{sec:1.2.4}
For many purposes it is useful to extend the system C of complex numbers by introduction of a symbol $\infty$ to represent infinity\index{Infinity}. Its connection with the finite numbers is established by setting $a + \infty = \infty + a = \infty$ for all finite a, and
\begin{equation*}
b \cdot \infty = \infty \cdot b = \infty
\end{equation*}
for all $b \neq 0$, including $b = \infty$. It is impossible, however, to define $\infty + \infty$ and $0 \cdot \infty$ without violating the laws of arithmetic. By special convention we shall nevertheless write $a/0 = \infty$ for $a \neq 0$ and $b/\infty = 0$ for $b \neq \infty$.

In the plane there is no room for a point corresponding to $\infty$, but we can of course introduce an “ideal” point which we call the point at infinity. The points in the plane together with the point at infinity form the extended complex plane. We agree that every straight line shall pass through the point at infinity. By contrast, no half plane shall contain the ideal point.

It is desirable to introduce a geometric model in which all points of the extended plane\index{Plane!extended} have a concrete representative. To this end we consider the unit sphere S whose equation in three-dimensional space is $x_{1}^{2} + x_{2}^{2} + x_{3}^{2} = 1$. With every point on S, except $(0,0,1)$, we can associate a complex number
\begin{equation*}
z = \frac {x _ {1} + i x _ {2}}{1 - x _ {3}},\tag{24}\label{eq:1.24}
\end{equation*}
and this correspondence is one to one. Indeed, from \hyperref[eq:1.24]{(24)} we obtain
\begin{equation*}
| z | ^ {2} = \frac {x _ {1} ^ {2} + x _ {2} ^ {2}}{(1 - x _ {3}) ^ {2}} = \frac {1 + x _ {3}}{1 - x _ {3}},
\end{equation*}
and hence
\begin{equation*}
x _ {3} = \frac {| z | ^ {2} - 1}{| z | ^ {2} + 1}.\tag{25}\label{eq:1.25}
\end{equation*}
Further computation yields
\begin{equation*}
\begin{array}{l} x _ {1} = \frac {z + \bar {z}}{1 + | z | ^ {2}} \\ x _ {2} = \frac {z - \bar {z}}{i (1 + | z | ^ {2})}. \end{array}\tag{26}\label{eq:1.26}
\end{equation*}
The correspondence can be completed by letting the point at infinity correspond to $(0,0,1)$, and we can thus regard the sphere as a representation of the extended plane or of the extended number system. We note that the hemisphere $x_{3} < 0$ corresponds to the disk $|z| < 1$ and the
% Original PDF physical page 35
\sourcepage{35}
hemisphere $x_{3} > 0$ to its outside $|z| > 1$. In function theory the sphere S is referred to as the Riemann sphere\index{Riemann sphere}.

If the complex plane is identified with the $(x_{1},x_{2})$ -plane with the $x_{1}$ - and $x_{2}$ -axis corresponding to the real and imaginary axis, respectively, the transformation \hyperref[eq:1.24]{(24)} takes on a simple geometric meaning. Writing $z = x + iy$ we can verify that
\begin{equation*}
x: y: - 1 = x _ {1}: x _ {2}: x _ {3} - 1,\tag{27}\label{eq:1.27}
\end{equation*}
and this means that the points $(x,y,0)$, $(x_{1},x_{2},x_{3})$, and $(0,0,1)$ are in a straight line. Hence the correspondence is a central projection from the center $(0,0,1)$ as shown in Fig.~\hyperref[fig:1-3]{1-3}. It is called a stereographic projection\index{Stereographic projection}. The context will make it clear whether the stereographic projection is regarded as a mapping from S to the extended complex plane, or vice versa.

In the spherical representation there is no simple interpretation of addition and multiplication. Its advantage lies in the fact that the point at infinity is no longer distinguished.

It is geometrically evident that the stereographic projection transforms every straight line in the z-plane into a circle on S which passes through the pole $(0,0,1)$, and the converse is also true. More generally, any circle on the sphere corresponds to a circle or straight line in the z-plane. To prove this we observe that a circle on the sphere lies in a plane $\alpha_{1}x_{1} + \alpha_{2}x_{2} + \alpha_{3}x_{3} = \alpha_{0}$, where we can assume that $\alpha_{1}^{2} + \alpha_{2}^{2} + \alpha_{3}^{2} = 1$ and $0 \leq \alpha_{0} < 1$. In terms of z and $\bar{z}$ this equation takes the form
\begin{equation*}
\alpha_ {1} (z + \bar {z}) - \alpha_ {2} i (z - \bar {z}) + \alpha_ {3} (| z | ^ {2} - 1) = \alpha_ {0} (| z | ^ {2} + 1)
\end{equation*}
or
\begin{equation*}
(\alpha_ {0} - \alpha_ {3}) (x ^ {2} + y ^ {2}) - 2 \alpha_ {1} x - 2 \alpha_ {2} y + \alpha_ {0} + \alpha_ {3} = 0.
\end{equation*}
For $\alpha_{0} \neq \alpha_{3}$ this is the equation of a circle, and for $\alpha_{0} = \alpha_{3}$ it represents a straight line. Conversely, the equation of any circle or straight line
\bookfigure{35}{1}{FIG. 1-3. Stereographic projection.}
% Original PDF physical page 36
\sourcepage{36}
can\index{Distance!spherical|(} be written in this form. The correspondence is consequently one to one.

It is easy to calculate the distance $d(z,z')$ between the stereographic projections of z and $z'$. If the points on the sphere are denoted by $(x_{1},x_{2},x_{3})$, $(x_{1}',x_{2}',x_{3}')$, we have first
\begin{equation*}
(x _ {1} - x _ {1} ^ {\prime}) ^ {2} + (x _ {2} - x _ {2} ^ {\prime}) ^ {2} + (x _ {3} - x _ {3} ^ {\prime}) ^ {2} = 2 - 2 \left(x _ {1} x _ {1} ^ {\prime} + x _ {2} x _ {2} ^ {\prime} + x _ {3} x _ {3} ^ {\prime}\right).
\end{equation*}
From (35) and (36) we obtain after a short computation
\begin{equation*}
\begin{array}{r l} x _ {1} x _ {1} ^ {\prime} + x _ {2} x _ {2} ^ {\prime} + x _ {3} x _ {3} ^ {\prime} & = \frac {(z + \bar {z}) (z ^ {\prime} + \bar {z} ^ {\prime}) - (z - \bar {z}) (z ^ {\prime} - \bar {z} ^ {\prime}) + (| z | ^ {2} - 1) (| z ^ {\prime} | ^ {2} - 1)}{(1 + | z | ^ {2}) (1 + | z ^ {\prime} | ^ {2})} \\ & = \frac {(1 + | z | ^ {2}) (1 + | z ^ {\prime} | ^ {2}) - 2 | z - z ^ {\prime} | ^ {2}}{(1 + | z | ^ {2}) (1 + | z ^ {\prime} | ^ {2})}. \end{array}
\end{equation*}
As a result we find that
\begin{equation*}
d (z, z ^ {\prime}) = \frac {2 | z - z ^ {\prime} |}{\sqrt {(1 + | z | ^ {2}) (1 + | z ^ {\prime} | ^ {2})}}.\tag{28}\label{eq:1.28}
\end{equation*}
For $z' = \infty$ the corresponding formula is
\begin{equation*}
d (z, \infty) = \frac {2}{\sqrt {1 + | z | ^ {2}}}.
\end{equation*}
\subsection*{EXERCISES}
\originalnumber{1}{exercise:1.2.4.8.1}
1. Show that $z$ and $z'$ correspond to diametrically opposite points on the Riemann sphere if and only if $z\bar{z}' = -1$.

\originalnumber{2}{exercise:1.2.4.8.2}
2. A cube has its vertices on the sphere S and its edges parallel to the coordinate axes. Find the stereographic projections of the vertices.

\originalnumber{3}{exercise:1.2.4.8.3}
3. Same problem for a regular tetrahedron in general position.

\originalnumber{4}{exercise:1.2.4.8.4}
4. Let $Z, Z'$ denote the stereographic projections of $z, z'$, and let $N$ be the north pole. Show that the triangles $NZZ'$ and $Nzz'$ are similar, and use this to derive \hyperref[eq:1.28]{(28)}.

\originalnumber{5}{exercise:1.2.4.8.5}
5. Find the radius of the spherical image of the circle in the plane whose center is a and radius R.

\index{Distance!spherical|)}